\documentclass[11pt]{article}
\usepackage[margin=1in]{geometry}
\usepackage{amsmath,amssymb,mathtools}
\usepackage{fontspec,unicode-math}
\setmainfont{Latin Modern Roman}
\setsansfont{Latin Modern Sans}
\setmathfont{Latin Modern Math}
\usepackage{xurl,hyperref}
\hypersetup{colorlinks=true,urlcolor=blue,linkcolor=blue}
\newcommand{\sref}[2]{\hyperlink{source:#1:#2}{[S#2]}}
\newcommand{\cataloguescope}[1]{\paragraph{Recorded mathematical question.}#1}
\begin{document}
\section{Growth exponents in the balanced city-growth regime}\label{statement}

\subsection{Definitions and mathematical statement}

\paragraph{Shared notation.}$\mathbb N=\{1,2,\ldots\}$, and $\mathbb Z,\mathbb Q,\mathbb R,\mathbb C$ denote the integers, rationals, reals and complex numbers. Any other conventions are specified locally.


At integer time $t\ge1$, cities have distinct positions $x_i\in[0,1]^2$ and populations $N_i(t)\ge1$ summing to $t$. Start with one city of population one at a uniform random position. Fix $c_0,\alpha,\beta>0$, and define
\[I(y,t)=c_0\max_i N_i(t)^\alpha\|y-x_i\|_2^{-\beta}.\]
An independent uniform point $U\in[0,1]^2$ is selected for immigrant $t+1$. With probability $1/(1+I(U,t))$ the immigrant founds a new city of population one at $U$; otherwise it joins a city attaining the maximum influence at $U$, increasing that city's population by one. Ties may be resolved arbitrarily on the zero-area boundaries. Write $M(t)$ for number of cities, $N_{(1)}(t)=\max_iN_i(t)$, and $N_{(\ge2)}(t)=t-N_{(1)}(t)$. The nondecreasing quantities $M(t)$ and $N_{(\ge2)}(t)$ have limits in $[0,\infty]$. Define $R(y,t)=\min_i\|y-x_i\|_2$, and let $S_i(t)$ be the influence cell where city $i$ attains the maximum defining $I(y,t)$; put $a_i(t)=\operatorname{area}S_i(t)$. For $Q=I,R$, define
\[\zeta^*[Q]=\inf\{u:\mathbb E\operatorname{area}\{y\in[0,1]^2:Q(y,t)\ge t^u\}\to0\},\]
\[\zeta_*[Q]=\sup\{u:\mathbb E\operatorname{area}\{y\in[0,1]^2:Q(y,t)\le t^u\}\to0\}.\]
For population, use the influence-cell weighted exponents
\[\zeta^*[N]=\inf\{u:\mathbb E\sum_i a_i(t)\mathbf1_{\{N_i(t)\ge t^u\}}\to0\},\qquad
\zeta_*[N]=\sup\{u:\mathbb E\sum_i a_i(t)\mathbf1_{\{N_i(t)\le t^u\}}\to0\}.\]
All limits in these definitions are as $t\to\infty$.
\paragraph{Statement recorded in the supplied source.}
For $0<\alpha<1$ and $\beta>2\alpha$, prove
\[\zeta_*[I]=\zeta^*[I]=\zeta_*[N]=\zeta^*[N]=\frac{\beta}{2-2\alpha+\beta},\qquad
\zeta_*[R]=\zeta^*[R]=\frac{\alpha-1}{2-2\alpha+\beta}.\] \sref{statement}{2}
\subsection{Short English statement}
Identify all six upper and lower influence, population and nearest-city distance exponents in the balanced growth regime.
\subsection{Sources}
\begin{description}
\item[\hypertarget{source:statement:1}{S1}] ULAM.AI and the UnsolvedMath Contributors, \emph{UnsolvedMath: A Curated Collection of Open Mathematics Problems}, catalogue entry. CC BY 4.0. \url{https://huggingface.co/datasets/ulamai/UnsolvedMath}.
\item[\hypertarget{source:statement:2}{S2}] David J. Aldous, technical notes on the spatial city-growth model, draft dated 25 April 2007, Section 1, model, p. 2; Section 3, exponent definitions and conjectured values, p. 15. \url{https://www.stat.berkeley.edu/~aldous/Research/OP/cities-notes.pdf}.
\end{description}
\label{end:statement}
\end{document}
